% C07 conceptual visual: branch cuts and monodromy.
\begin{figure}[htbp]
\centering
\begin{tikzpicture}[scale=1.0,>=Stealth]
  \draw[->] (-3.2,0) -- (3.2,0) node[right] {$\Re z$};
  \draw[->] (0,-2.2) -- (0,2.2) node[above] {$\Im z$};
  \fill (0,0) circle (1.5pt) node[below right] {$0$};
  \draw[line width=1.1pt] (0,0) -- (-3.0,0);
  \node[below] at (-2.1,0) {branch cut};
  \draw[->,line width=0.9pt] (1.35,0) arc (0:330:1.35);
  \node at (1.9,1.45) {$\arg z\mapsto \arg z+2\pi$};
\end{tikzpicture}
\caption{A logarithm becomes single-valued after a ray is removed. On the punctured
plane a loop winding once around the branch point changes a continued logarithm by
$2\pi i$; the slit prevents that loop from remaining inside the chosen branch domain.}
\label{fig:iv14-c07-branch-cut}
\end{figure}

\begin{table}[htbp]
\centering\small
\begin{tabular}{p{0.27\linewidth}p{0.61\linewidth}}
\toprule
Choice & Convention and consequence\\
\midrule
Branch domain & A simply connected domain in $\mathbb C^\times$ supports a holomorphic logarithm.\\
Principal argument & $\operatorname{Arg} z\in(-\pi,\pi]$ places the cut on the nonpositive real axis.\\
Root convention & $z^{1/n}=\exp(\operatorname{Log} z/n)$ is single-valued only after the logarithm branch is fixed.\\
Monodromy test & Continuation around a nontrivial winding detects whether a global branch can exist.\\
\bottomrule
\end{tabular}
\caption{Branch notation is meaningful only together with its domain and argument convention.}
\label{tab:iv14-c07-branch-conventions}
\end{table}
